% Part I -- Getting Started
% Chapter 1: Your First Python Program

\h{Your First Python Program}

Python lets you give a computer a clear sequence of instructions. Your goal in
this chapter is simple: install the necessary tools, run a short program, and
understand what appears on the screen. You do not need previous coding
experience.

\begin{NoteBox}
\textbf{Prerequisites.} A Windows or macOS computer, permission to install
software, and a folder where you can save your work.
\end{NoteBox}

\begin{tcolorbox}[title=Learning outcomes,colframe=orange,colback=white,arc=2mm]
After completing this chapter, you should be able to:
\begin{itemize}
    \item explain the difference between the editor, terminal, script, and REPL;
    \item verify that Python 3 is installed;
    \item create and run a Python file in Visual Studio Code;
    \item use \c{print()} to display text and calculation results; and
    \item respond to common setup errors without guessing.
\end{itemize}
\end{tcolorbox}

\hh{How to use this book}

This book is designed for independent study. Work through each chapter in
order. Type the examples instead of only reading them. Small typing mistakes
are normal, and correcting them is part of learning.

You will see five kinds of material:
\begin{itemize}
    \item \textbf{Explanation}: a concept described in plain language.
    \item \textbf{Runnable code}: Python you can type into a file and run.
    \item \textbf{Expected output}: what a correct run should display.
    \item \textbf{Stop and predict}: a question to answer before running code.
    \item \textbf{Practice}: an activity followed by hints and a solution.
\end{itemize}

Do not try to memorize every command. Focus on recognizing what each line asks
Python to do.

\hh{What Python is}

Python is a programming language. A Python program is a text file containing
instructions. Python reads those instructions and performs them in order.

\code{c01_first_print}{python}{
    print("Hello, Python!")
}{A single Python instruction}

The word \c{print} names an action. The parentheses contain the information the
action will use. The quotation marks identify text.

\begin{NoteBox}
\textbf{Stop and predict.} What would you change to display
\c{Hello, ARC500!}?
\end{NoteBox}

\hh{Install Python 3}

Python 3 is the form of the language used in this book.

\hhh*{Windows}
\begin{enumerate}
    \item Visit \url{https://www.python.org/downloads/}.
    \item Download the current Python 3 installer.
    \item Start the installer.
    \item Select \textbf{Add Python to PATH} before choosing Install.
\end{enumerate}

\hhh*{macOS}
\begin{enumerate}
    \item Visit \url{https://www.python.org/downloads/}.
    \item Download and run the current macOS installer.
    \item Complete the installation using the default options.
\end{enumerate}

\hh{Verify the installation}

Open a terminal and enter one of the following commands. Run the commands one
at a time. You only need one of them to work.

\code{c01_verify_python}{bash}{
    python --version
    py --version
    python3 --version
}{Terminal commands for checking the Python version}

A successful response will be similar to:

\code{c01_version_output}{text}{
    Python 3.x.x
}{Example version output; the numbers may differ}

These are terminal commands. Do not type them into a Python file.

\hh{Install Visual Studio Code}

Visual Studio Code, usually called VS Code, is the editor used in this book.
The editor is where you write and save Python instructions.

\begin{enumerate}
    \item Download VS Code from \url{https://code.visualstudio.com/}.
    \item Install and open it.
    \item Select the Extensions icon on the left side.
    \item Search for \textbf{Python} and install the extension published by Microsoft.
    \item Open the Command Palette with \texttt{Ctrl+Shift+P} on Windows or
          \texttt{Cmd+Shift+P} on macOS.
    \item Choose \textbf{Python: Select Interpreter}, then select Python 3.
\end{enumerate}

\hh{Create a study folder}

Create a folder named \texttt{arc500-python}. Inside it, create a folder named
\texttt{chapter01}. Keeping each chapter in its own folder makes your files
easier to find.

In VS Code, choose \textbf{File $\rightarrow$ Open Folder} and open
\texttt{chapter01}. The folder name should appear in the Explorer panel.

\hh{Understand the four working areas}

\begin{itemize}
    \item \textbf{Editor}: the area in VS Code where you write and save code.
    \item \textbf{Terminal}: a text-based area where you enter system commands.
    \item \textbf{Script}: a saved Python file whose name ends in \texttt{.py}.
    \item \textbf{REPL}: an interactive Python prompt that runs one instruction at a time.
\end{itemize}

The terminal and editor are not interchangeable. A command such as
\texttt{python hello.py} belongs in the terminal. An instruction such as
\c{print("Hello")} belongs in a Python file or the REPL.

\hh{Write and run your first script}

In the VS Code Explorer, create a file named \texttt{hello.py}. Type and save:

\code{c01_hello_script}{python}{
    print("Hello, ARC500!")
    print("This is my first Python script.")
}{The complete hello.py script}

Open the VS Code terminal with \textbf{Terminal $\rightarrow$ New Terminal}.
Run the script with the command that works on your computer:

\code{c01_run_script}{bash}{
    python hello.py
}{Run a Python script from the terminal}

On some Windows computers, use \texttt{py hello.py}. On some macOS computers,
use \texttt{python3 hello.py}.

Expected output:

\code{c01_hello_output}{text}{
    Hello, ARC500!
    This is my first Python script.
}{Expected output from hello.py}

Python runs the first instruction and then the second. The output appears in
the same order as the instructions.

\begin{NoteBox}
\textbf{Stop and predict.} If you exchange the two lines in the file, what will
appear first when you run it again?
\end{NoteBox}

\hh{Display a calculation}

The \c{print()} function can display a calculation as well as text.

\code{c01_print_calculations}{python}{
    print(6 + 4)
    print(6 * 4)
    print(24 / 6)
}{Display three calculation results}

Expected output:

\code{c01_calculation_output}{text}{
    10
    24
    4.0
}{Output from the calculations}

Python uses \c{*} for multiplication and \c{/} for division. Chapter 2
explains numbers and arithmetic in more detail.

\hh{Use the interactive prompt}

The REPL is useful when you want to try one instruction without creating a
file. In the terminal, enter \texttt{python}, \texttt{py}, or
\texttt{python3}. A prompt similar to \c{>>>} will appear.

\code{c01_repl_session}{text}{
    >>> print("Testing one instruction")
    Testing one instruction
    >>> 8 * 5
    40
}{A short REPL session}

The \c{>>>} symbol is the prompt. Do not include it in a Python script. Enter
\c{exit()} to leave the REPL.

\hh{Read an error message calmly}

Suppose you accidentally omit a closing quotation mark:

\code{c01_syntax_error}{python}{
    print("Hello)
}{An intentionally incorrect instruction}

Python reports a \c{SyntaxError}. The last line of the message gives the error
type. Look near the indicated line for a missing quotation mark, parenthesis,
or other symbol.

\code{c01_syntax_fixed}{python}{
    print("Hello")
}{The corrected instruction}

\hh{Troubleshooting setup problems}

\begin{itemize}
    \item \textbf{Command not found}: try \texttt{py} or \texttt{python3}. If none
          works, reinstall Python and ensure it is added to PATH.
    \item \textbf{The file does not run}: confirm that the terminal is open in
          the folder containing the file.
    \item \textbf{Old output appears}: save the file before running it.
    \item \textbf{The wrong version appears}: use
          \textbf{Python: Select Interpreter} in VS Code.
    \item \textbf{Nothing appears}: confirm that the file contains \c{print()}
          and that you ran the correct filename.
\end{itemize}

\hh{Practice ladder}

\hhh{Level 0 -- Read and predict}

Write down the expected output before running this code:

\code{c01_practice_level0}{python}{
    print("Start")
    print(3 + 2)
    print("Finish")
}{Predict the output before running the code}

\hhh{Level 1 -- Modify}

Copy the code into \texttt{practice1.py}. Replace the text with your own
subject and learning goal.

\code{c01_practice_level1}{python}{
    print("Subject: Python")
    print("Goal: Run a program successfully")
}{Modify two text values}

\hhh{Level 2 -- Complete}

Write three \c{print()} instructions that produce:

\code{c01_practice_level2_output}{text}{
    ARC500
    20
    Complete
}{Target output for Level 2}

\hhh{Level 3 -- Apply}

Create \texttt{practice3.py}. It should display your name, the text
\texttt{Chapter 1 complete}, and the result of \c{9 * 7}.

\textbf{Hints}
\begin{itemize}
    \item \textbf{Hint 1}: Use one \c{print()} instruction for each line.
    \item \textbf{Hint 2}: Text needs quotation marks; a calculation does not.
    \item \textbf{Hint 3}: Your file should contain three instructions.
\end{itemize}

\hhh{Level 4 -- Challenge}

Use four \c{print()} instructions to produce:

\code{c01_practice_level4_output}{text}{
    Python Study Record
    -------------------
    Chapter: 1
    Status: Complete
}{Target output for the optional challenge}

\hh{Practice solutions}

\textbf{Level 0}

\code{c01_solution_level0}{text}{
    Start
    5
    Finish
}{Level 0 solution}

\textbf{Level 2}

\code{c01_solution_level2}{python}{
    print("ARC500")
    print(12 + 8)
    print("Complete")
}{One solution to Level 2}

\textbf{Level 3}

\code{c01_solution_level3}{python}{
    print("Your Name")
    print("Chapter 1 complete")
    print(9 * 7)
}{One solution to Level 3}

\textbf{Level 4}

\code{c01_solution_level4}{python}{
    print("Python Study Record")
    print("-------------------")
    print("Chapter: 1")
    print("Status: Complete")
}{Solution to Level 4}

\hh{Chapter checkpoint}

Answer without looking back, then check your answers.

\begin{enumerate}
    \item What filename extension identifies a Python script?
    \item Where do you write and save Python code in VS Code?
    \item Where do you enter \texttt{python hello.py}?
    \item What does \c{print()} do?
    \item What symbols identify text in this chapter's examples?
    \item What multiplication symbol does Python use?
    \item What is the purpose of the REPL?
    \item Should \c{>>>} be copied into a Python script?
    \item What should you check first when old output appears?
    \item Name one common cause of a \c{SyntaxError}.
\end{enumerate}

\hh{Checkpoint answers}

\begin{enumerate}
    \item \texttt{.py}.
    \item In the editor.
    \item In the terminal.
    \item It displays information.
    \item Quotation marks.
    \item An asterisk: \c{*}.
    \item It runs one instruction or expression at a time.
    \item No. It is the interactive prompt.
    \item Check that the file was saved and that you ran the correct file.
    \item A missing quotation mark or parenthesis is one possible cause.
\end{enumerate}

\begin{tcolorbox}[title=Chapter summary,colframe=orange,colback=white,arc=2mm]
You installed Python and VS Code, created and ran a Python file, displayed text
and calculations, used the REPL, and corrected a basic syntax mistake. Chapter
2 introduces values, variables, and calculations.
\end{tcolorbox}

